\begin{thebibliography}{99}

\bibitem[Abramowitz and Stegun(1964)]{AbramowitzStegun1964} Abramowitz, M.; Stegun, I.A. \textit{Handbook of Mathematical Functions with Formulas, Graphs, and Mathematical Tables}; National Bureau of Standards: Washington, DC, 1964.

\bibitem[Ahn et~al.(2012)]{Ahn2012} Ahn, S.; Kim, J.H.T.; Ramaswami, V. A new class of models for heavy tailed distributions in finance and insurance risk. \textit{Insur. Math. Econ.} 2012, \textit{51}, 43--52.

\bibitem[Akaike(1974)]{Akaike1974} Akaike, H. A new look at the statistical model identification. \textit{IEEE Trans. Autom. Control} 1974, \textit{19}, 716--723.

\bibitem[Albrecher and Bladt(2019)]{AlbrecherBladt2019} Albrecher, H.; Bladt, M. Inhomogeneous phase-type distributions and heavy tails. \textit{J. Appl. Probab.} 2019, \textit{56}, 1044--1064.

\bibitem[Asmussen et~al.(1996)]{Asmussen1996} Asmussen, S.; Nerman, O.; Olsson, M. Fitting phase-type distributions via the EM algorithm. \textit{Scand. J. Stat.} 1996, \textit{23}, 419--441.

\bibitem[Balkema and de Haan(1974)]{BalkemaDeHaan1974} Balkema, A.A.; de Haan, L. Residual life time at great age. \textit{Ann. Probab.} 1974, \textit{2}, 792--804.

\bibitem[Beirlant et~al.(1999)]{Beirlant1999} Beirlant, J.; Dierckx, G.; Goegebeur, Y.; Matthys, G. Tail index estimation and an exponential regression model. \textit{Extremes} 1999, \textit{2}, 177--200.

\bibitem[Beirlant et~al.(2004)]{Beirlant2004} Beirlant, J.; Goegebeur, Y.; Segers, J.; Teugels, J. \textit{Statistics of Extremes: Theory and Applications}; Wiley: Chichester, 2004.

\bibitem[Bladt and Nielsen(2017)]{BladtNielsen2017} Bladt, M.; Nielsen, B.F. \textit{Matrix-Exponential Distributions in Applied Probability}; Springer: New York, 2017.

\bibitem[Coles(2001)]{Coles2001} Coles, S. \textit{An Introduction to Statistical Modeling of Extreme Values}; Springer: London, 2001.

\bibitem[Dickson and Hipp(2000)]{DicksonHipp2000} Dickson, D.C.M.; Hipp, C. Ruin problems for phase-type(2) risk processes. \textit{Scand. Actuar. J.} 2000, \textit{2000}, 147--167.

\bibitem[Drees et~al.(2000)]{Drees2000} Drees, H.; de Haan, L.; Resnick, S. How to make a Hill plot. \textit{Ann. Stat.} 2000, \textit{28}, 254--274.

\bibitem[Embrechts et~al.(1997)]{Embrechts1997} Embrechts, P.; Kl{\"u}ppelberg, C.; Mikosch, T. \textit{Modelling Extremal Events for Insurance and Finance}; Springer: Berlin, 1997.

\bibitem[Feuerverger and Hall(1999)]{FeuervergerHall1999} Feuerverger, A.; Hall, P. Estimating a tail exponent by modelling departure from a Pareto distribution. \textit{Ann. Stat.} 1999, \textit{27}, 760--781.

\bibitem[Fraga Alves et~al.(2003)]{FragaAlves2003} Fraga Alves, M.I.; Gomes, M.I.; de Haan, L. A new class of semi-parametric estimators of the second order parameter. \textit{Port. Math.} 2003, \textit{60}, 193--213.

\bibitem[Geluk et~al.(1997)]{Geluk1997} Geluk, J.; de Haan, L.; Resnick, S.; St{\u{a}}ric{\u{a}}, C. Second-order regular variation, convolution and the central limit theorem. \textit{Stoch. Process. Appl.} 1997, \textit{69}, 139--159.

\bibitem[Goegebeur et~al.(2010)]{Goegebeur2010} Goegebeur, Y.; Beirlant, J.; de Wet, T. Kernel estimators for the second order parameter in extreme value statistics. \textit{J. Stat. Plan. Inference} 2010, \textit{140}, 2632--2652.

\bibitem[Gomes and Guillou(2015)]{GomesGuillou2015} Gomes, M.I.; Guillou, A. Extreme value theory and statistics of univariate extremes: a review. \textit{Int. Stat. Rev.} 2015, \textit{83}, 263--292.

\bibitem[De Haan and Ferreira(2006)]{DeHaanFerreira2006} de Haan, L.; Ferreira, A. \textit{Extreme Value Theory: An Introduction}; Springer: New York, 2006.

\bibitem[Hall(1982)]{Hall1982} Hall, P. On some simple estimates of an exponent of regular variation. \textit{J. R. Stat. Soc. Ser. B} 1982, \textit{44}, 37--42.

\bibitem[Hill(1975)]{Hill1975} Hill, B.M. A simple general approach to inference about the tail of a distribution. \textit{Ann. Stat.} 1975, \textit{3}, 1163--1174.

\bibitem[Kim and Kim(2015)]{KimKim2015} Kim, J.H.T.; Kim, J. A parametric alternative to the Hill estimator for heavy-tailed distributions. \textit{J. Bank. Finance} 2015, \textit{54}, 60--71.

\bibitem[McNeil(1997)]{McNeil1997} McNeil, A.J. Estimating the tails of loss severity distributions using extreme value theory. \textit{ASTIN Bull.} 1997, \textit{27}, 117--137.

\bibitem[Neuts(1981)]{Neuts1981} Neuts, M.F. \textit{Matrix-Geometric Solutions in Stochastic Models: An Algorithmic Approach}; Johns Hopkins University Press: Baltimore, 1981.

\bibitem[Pickands(1975)]{Pickands1975} Pickands, J. Statistical inference using extreme order statistics. \textit{Ann. Stat.} 1975, \textit{3}, 119--131.

\bibitem[Resnick(1997)]{Resnick1997} Resnick, S.I. Discussion of the Danish data on large fire insurance losses. \textit{ASTIN Bull.} 1997, \textit{27}, 139--151.

\bibitem[Reynkens et~al.(2017)]{Reynkens2017} Reynkens, T.; Verbelen, R.; Beirlant, J.; Antonio, K. Modelling censored losses using splicing: a global fit strategy with mixed Erlang and extreme value distributions. \textit{Insur. Math. Econ.} 2017, \textit{77}, 65--77.

\bibitem[Rotnitzky et~al.(2000)]{Rotnitzky2000} Rotnitzky, A.; Cox, D.R.; Bottai, M.; Robins, J. Likelihood-based inference with singular information matrix. \textit{Bernoulli} 2000, \textit{6}, 243--284.

\bibitem[Scarrott and MacDonald(2012)]{ScarrottMacDonald2012} Scarrott, C.; MacDonald, A. A review of extreme value threshold estimation and uncertainty quantification. \textit{REVSTAT Stat. J.} 2012, \textit{10}, 33--60.

\bibitem[Schwarz(1978)]{Schwarz1978} Schwarz, G. Estimating the dimension of a model. \textit{Ann. Stat.} 1978, \textit{6}, 461--464.

\bibitem[Self and Liang(1987)]{SelfLiang1987} Self, S.G.; Liang, K.-Y. Asymptotic properties of maximum likelihood estimators and likelihood ratio tests under nonstandard conditions. \textit{J. Am. Stat. Assoc.} 1987, \textit{82}, 605--610.

\bibitem[Smith(1987)]{Smith1987} Smith, R.L. Estimating tails of probability distributions. \textit{Ann. Stat.} 1987, \textit{15}, 1174--1207.

\bibitem[Weissman(1978)]{Weissman1978} Weissman, I. Estimation of parameters and large quantiles based on the $k$ largest observations. \textit{J. Am. Stat. Assoc.} 1978, \textit{73}, 812--815.

\end{thebibliography}
